\documentclass[11pt,a4paper]{article}

\usepackage[top=2.5cm, bottom=2.5cm, left=2.5cm, right=2.5cm]{geometry}
\usepackage{amsmath, amssymb}
\usepackage{graphicx}
\usepackage{booktabs}
\usepackage{array}
\usepackage{enumitem}
\usepackage{longtable}
\usepackage{xcolor}
\usepackage[hidelinks]{hyperref}
\usepackage{xurl}
\usepackage{microtype}

\setlist[enumerate]{itemsep=4pt, topsep=4pt, leftmargin=2em}
\setlist[itemize]{itemsep=3pt, topsep=4pt}
\setlength{\emergencystretch}{3em}

\title{\textbf{Proyek Akhir}\\Internet of Things (MK38)}
\author{Program Studi Informatika, Universitas Jambi}
\date{}

\begin{document}

\maketitle
\vspace{-2em}

\begin{center}
\begin{tabular}{ll}
\toprule
\textbf{Item} & \textbf{Keterangan} \\
\midrule
Mata kuliah & Internet of Things (MK38), 3 SKS, Semester 5 \\
Bentuk & Proyek kelompok (\textbf{6 kelompok $\times$ 5 mahasiswa}) \\
Dosen Pengampu & Edi Saputra, S.T., M.Sc. \\
 & Akhiyar Waladi, S.Komp., M.Kom. \\
Bobot & Milestone Proyek 10\% + Proyek Akhir 40\% \\
Alur kerja & Simulasi \texttt{Wokwi} lebih dahulu, baru implementasi \\
 & perangkat keras (hardware) \\
Pengerjaan & Pertemuan 11 sampai 15 (bimbingan dan milestone) \\
Presentasi & Pertemuan 15 (demonstrasi langsung) \\
\bottomrule
\end{tabular}
\end{center}

\section{Deskripsi Umum}

Proyek akhir adalah puncak mata kuliah Internet of Things. Setiap kelompok
membangun \textbf{satu sistem IoT end-to-end} untuk satu kasus nyata: mulai
dari perancangan (blok diagram, skematik \texttt{Fritzing}, desain 3D),
implementasi perangkat keras berbasis \textbf{ESP32}, komunikasi nirkabel ke
cloud menggunakan \textbf{MQTT} dan/atau \textbf{REST API}, sampai dashboard
pemantauan dan kendali jarak jauh.

Seluruh tahap pengembangan mengikuti alur kerja \emph{simulate first}: setiap
rangkaian \textbf{wajib disimulasikan di \texttt{Wokwi}}
(\url{https://wokwi.com}) sebelum dirakit pada perangkat keras, sehingga
kesalahan wiring dan program dapat dideteksi lebih awal tanpa merusak
komponen.

\section{Tujuan}

\begin{enumerate}
  \item Menerapkan siklus lengkap sistem IoT: sensor $\rightarrow$
        mikrokontroler $\rightarrow$ komunikasi nirkabel $\rightarrow$ cloud
        $\rightarrow$ aktuator/dashboard.
  \item Melatih kerja tim dalam kelompok besar (5 orang) dengan pembagian
        peran yang jelas (perangkat keras, firmware, cloud/dashboard,
        dokumentasi, pengujian).
  \item Menghasilkan prototype yang berfungsi, terdokumentasi, dan dapat
        didemonstrasikan secara langsung.
\end{enumerate}

\section{Ketentuan Kelompok dan Topik}

\begin{enumerate}
  \item Kelas dibagi menjadi \textbf{6 kelompok}, setiap kelompok
        \textbf{5 mahasiswa}. Pembagian dilakukan pada pertemuan 11.
  \item Setiap kelompok memilih \textbf{satu topik berbeda} dari enam topik
        pada Bagian~\ref{sec:topik} (satu topik hanya boleh diambil satu
        kelompok).
  \item Usulan topik lain diperbolehkan dengan tingkat kesulitan setara dan
        persetujuan dosen.
  \item Semua topik telah diseimbangkan bobot kerjanya: jumlah sensor,
        aktuator, protokol komunikasi, dan keluaran yang wajib dibuat
        \textbf{setara} (lihat matriks keseimbangan pada
        Bagian~\ref{sec:balance}).
\end{enumerate}

\section{Topik Proyek}\label{sec:topik}

Setiap topik memiliki: masalah yang dipecahkan, daftar komponen, protokol
komunikasi, dan metrik keberhasilan. Fitur disusun \textbf{bertingkat}:

\begin{itemize}
  \item \textbf{Tingkat 1 --- Dasar (wajib):} inti sistem yang harus berjalan;
        menjadi batas minimal kelulusan fungsi.
  \item \textbf{Tingkat 2 --- Menengah (pilih minimal 2):} fitur integrasi
        cloud/dashboard yang menaikkan nilai fungsionalitas.
  \item \textbf{Tingkat 3 --- Lanjutan (bonus):} fitur pembeda; dikerjakan
        bila waktu dan kemampuan mendukung.
\end{itemize}

Struktur bertingkat ini membuat \textbf{upaya tetap terlihat dan dinilai}
meskipun kelompok tidak menuntaskan seluruh tingkat: capaian Tiap tingkat
didokumentasikan pada progress report, dan penilaian menghitung apa yang
berhasil dibangun dengan baik --- bukan hanya apa yang gagal dicapai.

\subsection*{Topik 1 --- Smart Parking dengan Payment Gateway}

\begin{itemize}
  \item \textbf{Masalah:} Pengemudi sulit mengetahui slot parkir kosong dan
        proses bayar manual lambat. Sistem mendeteksi slot secara otomatis,
        membuka palang saat masuk/keluar, dan menghitung tarif dengan
        pembayaran digital.
  \item \textbf{Tingkat 1 --- Dasar (wajib):} (1) deteksi 4--6 slot parkir;
        (2) palang servo otomatis di pintu masuk/keluar; (3) tampilan sisa
        slot pada LCD/OLED; (4) publish status slot ke broker MQTT.
  \item \textbf{Tingkat 2 --- Menengah (pilih minimal 2):} verifikasi
        kartu RFID RC522 untuk akses member; riwayat masuk/keluar tersimpan ke
        database; indikator LED merah/hijau per slot; mode offline (palang
        tetap bekerja tanpa internet).
  \item \textbf{Tingkat 3 --- Lanjutan (bonus):} simulasi pembayaran
        \emph{payment gateway} (sandbox Midtrans/Xendit atau QR mock) yang
        membuka palang keluar setelah lunas; pengenalan plat nomor sederhana
        dengan ESP32-CAM; analitik okupansi harian di dashboard.
  \item \textbf{Komponen:} ESP32 ($\times$2: gerbang + slot), sensor
        ultrasonik/IR ($\times$4--6), servo ($\times$2), LCD/OLED I2C,
        RFID RC522 (opsional), buzzer.
  \item \textbf{Protokol:} MQTT (status slot, kendali palang) + REST API
        (verifikasi pembayaran); dashboard web/Node-RED.
  \item \textbf{Metrik:} akurasi deteksi slot ($\geq$ 95\%), waktu respon
        buka palang ($\leq$ 2 detik), alur bayar-palang berjalan utuh.
  \item \textbf{Sisi menyenangkan:} maket parkir mini dengan mobil-mobilan;
        scan QR palsu untuk ``membayar''.
\end{itemize}

\subsection*{Topik 2 --- Smart Greenhouse Otomatis}

\begin{itemize}
  \item \textbf{Masalah:} Tanaman greenhouse butuh suhu, kelembapan tanah, dan
        penyiraman terkontrol; petani tidak bisa memantau 24 jam. Sistem
        menyiram dan menjentikkan ventilasi otomatis sesuai ambang batas.
  \item \textbf{Tingkat 1 --- Dasar (wajib):} (1) monitoring suhu/kelembapan
        udara (DHT22); (2) kelembapan tanah (soil moisture); (3) pompa siram
        otomatis lewat relay saat tanah kering; (4) kipas ventilasi otomatis
        saat panas.
  \item \textbf{Tingkat 2 --- Menengah (pilih minimal 2):} dashboard MQTT
        (Node-RED/Grafana) dengan grafik real-time; notifikasi Telegram bila
        tangki air kritis (level HC-SR04); tampilan status lokal OLED;
        pencahayaan tumbuh otomatis mengikuti LDR.
  \item \textbf{Tingkat 3 --- Lanjutan (bonus):} data historis tersimpan ke
        InfluxDB/ThingSpeak; jadwal penyiraman berbasis waktu via dashboard;
        node sensor ESP-NOW terpisah untuk beda rak tanaman; rekomendasi
        penyiraman berbasis tren kelembapan.
  \item \textbf{Komponen:} ESP32, DHT22, soil moisture capacitive, LDR,
        HC-SR04 (level tangki), relay ($\times$2), pompa mini 5V, kipas DC,
        OLED SSD1306, tanaman pot kecil.
  \item \textbf{Protokol:} MQTT ke broker Mosquitto/HiveMQ; data historis ke
        Node-RED/InfluxDB; notifikasi via Telegram Bot API (REST).
  \item \textbf{Metrik:} kelembapan tanah terjaga pada rentang target
        ($\geq$ 90\% waktu), notifikasi Telegram terkirim 100\%, grafik
        historis terekam.
  \item \textbf{Sisi menyenangkan:} benar-benar merawat tanaman hidup selama
        proyek; demo menyiram otomatis di depan kelas.
\end{itemize}

\subsection*{Topik 3 --- Smart Home Terintegrasi Aplikasi Mobile/Web}

\begin{itemize}
  \item \textbf{Masalah:} Perabot rumah masih dikendalikan manual padahal
        penghuni ingin kontrol dan pemantauan energi dari ponsel, termasuk
        mode otomatis ketika ruangan kosong.
  \item \textbf{Tingkat 1 --- Dasar (wajib):} (1) kendali 3--4 peralatan
        (relay lampu/kipas); (2) deteksi keberadaan PIR untuk mode otomatis;
        (3) monitoring suhu/kelembapan ruangan; (4) aplikasi web/mobile
        dashboard dengan tombol kendali real-time.
  \item \textbf{Tingkat 2 --- Menengah (pilih minimal 2):} monitoring daya
        listrik per perangkat (PZEM-004T); scene otomatis (mode malam / mode
        pergi); log aktivitas kendali tersimpan; notifikasi anomali (PIR
        mendeteksi orang saat mode pergi).
  \item \textbf{Tingkat 3 --- Lanjutan (bonus):} integrasi voice assistant
        (Google Assistant/IFTTT); ESP32-CAM doorbell dengan deteksi wajah;
        update firmware over-the-air (OTA); analitik konsumsi energi mingguan.
  \item \textbf{Komponen:} ESP32 ($\times$1--2), relay 4-channel, PIR HC-SR501,
        DHT22, PZEM-004T (opsional, pengukur daya), buzzer.
  \item \textbf{Protokol:} REST API (kendali dari aplikasi) + MQTT (telemetri);
        dashboard web (HTML/Node-RED/Blynk).
  \item \textbf{Metrik:} latensi kendali aplikasi-perangkat ($\leq$ 1 detik),
        mode otomatis aktif tepat saat ruangan kosong/logis, semua relay
        dapat dikendalikan dari luar jaringan lokal.
  \item \textbf{Sisi menyenangkan:} mengendalikan ``rumah'' dari HP teman di
        kursi sebelah; mode malam otomatis yang bisa dipamerkan.
\end{itemize}

\subsection*{Topik 4 --- Adaptive Traffic Light Persimpungan Cerdas}

\begin{itemize}
  \item \textbf{Masalah:} Lampu lalu lintas durasi tetap menyia-nyiakan waktu
        di jalanan sepi dan memadatkan jalanan ramai. Sistem memberi hijau
        dinamis sesuai antrian tiap lengan.
  \item \textbf{Tingkat 1 --- Dasar (wajib):} (1) model persimpangan 4
        lengan; (2) sensor antrean per lengan (IR/PIR near-far); (3) lampu
        R/Y/G per lengan dengan siklus aman (all-red delay); (4) algoritma
        prioritas antrean dasar.
  \item \textbf{Tingkat 2 --- Menengah (pilih minimal 2):} dashboard MQTT
        menampilkan antrean dan fase lampu real-time; mode preemption saat
        antrean $\geq$ ambang; countdown waktu hijau pada LCD per lengan;
        logging statistik antrean per jam.
  \item \textbf{Tingkat 3 --- Lanjutan (bonus):} mode darurat ambulans
        (RFID/tombol) dengan prioritas tertinggi; sinkronisasi green-wave dua
        simpang (master-slave ESP-NOW); penghitungan kendaraan ESP32-CAM;
        prediksi densitas berbasis data historis.
  \item \textbf{Komponen:} ESP32 ($\times$1--2 master-slave), sensor IR/PIR
        ($\times$8), LED R/Y/G ($\times$12), LCD antrean, buzzer (mode
        darurat).
  \item \textbf{Protokol:} MQTT (skor antrean, fase lampu) + web dashboard;
        opsional ESP-NOW antar controller.
  \item \textbf{Metrik:} penurunan rata-rata waktu tunggu simulasi vs.\ timer
        tetap ($\geq$ 20\%), fase transisi selalu melalui all-red, prioritas
        ambulans bekerja.
  \item \textbf{Sisi menyenangkan:} maket persimpangan dengan mobil mainan;
        mengetahui bahwa algoritma kalian ``menyelamatkan'' ambulans.
\end{itemize}

\subsection*{Topik 5 --- Sistem Peringatan Dini Banjir / Flood Early Warning System}

\begin{itemize}
  \item \textbf{Masalah:} Banjir kilat menerpa tanpa peringatan. Sistem
        memantau ketinggian air dan intensitas hujan, lalu memberi sirene dan
        notifikasi bertingkat ke warga.
  \item \textbf{Tingkat 1 --- Dasar (wajib):} (1) ketinggian air via
        ultrasonik (3 level: aman\slash waspada\slash bahaya); (2) sensor
        hujan (rain drop module); (3) sirene buzzer + lampu status otomatis
        per level; (4) notifikasi Telegram Bot dengan perintah /status.
  \item \textbf{Tingkat 2 --- Menengah (pilih minimal 2):} dashboard MQTT
        dengan grafik riwayat level air; pola sirene berbeda per level
        bahaya; broadcast Telegram berkala tiap jam; relay + pompa mini
        sebagai demo drainase otomatis.
  \item \textbf{Tingkat 3 --- Lanjutan (bonus):} node sensor terpisah
        (ESP-NOW/LoRa) untuk titik pantau jauh; sensor debit aliran air;
        integrasi data cuaca publik via API; maket menara sirine mini.
  \item \textbf{Komponen:} ESP32, HC-SR04 (non-kontak, anti korosi), rain
        sensor, buzzer + LED RGB, relay + pompa mini (demo drainase), OLED.
  \item \textbf{Protokol:} REST API (Telegram Bot API) + MQTT (telemetri ke
        dashboard); debounce/filter agar tidak alarm palsu.
  \item \textbf{Metrik:} klasifikasi level akurat ($\geq$ 95\% uji), notifikasi
        terkirim $\leq$ 5 detik setelah level berubah, tanpa false alarm pada
        uji 30 menit.
  \item \textbf{Sisi menyenangkan:} membuat ``hujan'' dengan spray di atas
        sensor; grup kelas menerima alert Telegram sungguhan saat demo.
\end{itemize}

\subsection*{Topik 6 --- Smart Health Monitoring Portabel}

\begin{itemize}
  \item \textbf{Masalah:} Pemeriksaan detak jantung dan SpO$_2$ butuh alat
        mahal; pasien rawat inap/puskesmas desa perlu pemantauan jarak jauh
        dengan peringatan otomatis ke petugas.
  \item \textbf{Tingkat 1 --- Dasar (wajib):} (1) BPM dan SpO$_2$ via
        MAX30100\slash 30102; (2) suhu tubuh DS18B20\slash MLX90614;
        (3) tampilan lokal OLED; (4) publish data ke broker MQTT.
  \item \textbf{Tingkat 2 --- Menengah (pilih minimal 2):} dashboard grafik
        tren BPM/SpO$_2$ real-time; alarm otomatis (buzzer + Telegram
        petugas) bila nilai keluar ambang normal; profil multi-pasien dengan
        pemilihan via tombol; ekspor data pengukuran ke CSV.
  \item \textbf{Tingkat 3 --- Lanjutan (bonus):} eskalasi bertingkat
        (petugas $\rightarrow$ dokter) bila tidak ditanggapi; deteksi jatuh
        MPU6050 untuk pasien rawat; casing wearable hasil 3D print;
        mode hemat daya deep-sleep dengan bangun berkala.
  \item \textbf{Komponen:} ESP32, MAX30100/30102, DS18B20 waterproof, OLED
        SSD1306, buzzer, button (mode ukur).
  \item \textbf{Protokol:} MQTT (BPM, SpO$_2$, suhu) + REST (Telegram alert);
        dashboard Node-RED/Grafana.
  \item \textbf{Metrik:} BPM terbaca konsisten vs.\ oximeter referensi
        (selisih $\leq$ 5 bpm), alarm terpicu tepat pada uji ambang,
        data terekam kontinu $\geq$ 30 menit.
  \item \textbf{Sisi menyenangkan:} mengukur detak jantung anggota kelompok
        sebelum dan sesudah lari keliling gedung; data langsung muncul di
        layar dosen.
\end{itemize}

\subsection*{Matriks Keseimbangan Bobot Antar Topik}\label{sec:balance}

Semua topik dirancang setara pada lima dimensi berikut:

\begin{center}
\small
\begin{tabular}{@{}lccccc@{}}
\toprule
\textbf{Topik} & \textbf{Sensor} & \textbf{Aktuator} & \textbf{Komunikasi} &
\textbf{Dashboard} & \textbf{Tingkat 2} \\
\midrule
1 Parking        & 2 jenis & 3 jenis & MQTT + REST & ya & 4 pilihan \\
2 Greenhouse     & 4 jenis & 3 jenis & MQTT + REST & ya & 4 pilihan \\
3 Smart Home     & 2--3 jenis & relay $\times$4 & MQTT + REST & ya & 4 pilihan \\
4 Traffic Light  & 8 unit IR/PIR & LED 12 + LCD & MQTT & ya & 4 pilihan \\
5 Flood Warning  & 2 jenis & buzzer + relay & REST + MQTT & ya & 4 pilihan \\
6 Health Monitor & 2 jenis & buzzer + OLED & MQTT + REST & ya & 4 pilihan \\
\bottomrule
\end{tabular}
\end{center}

Catatan keseimbangan: Topik 4 memiliki jumlah unit sensor terbanyak namun
tiap unit sederhana (digital read); Topik 2 dan 6 memiliki sensor kalibrasi
yang sedikit lebih sensitif namun jumlahnya lebih sedikit; Topik 1 dan 3
memuat kompleksitas integrasi layanan eksternal (payment sandbox/aplikasi)
sebagai penggani kompleksitas sensor. Estimasi beban kerja tiap topik
setara $\pm$ 40--50 jam-jam-orang.

\section{Alur Pengerjaan (Milestone)}

\begin{center}
\small
\setlength{\tabcolsep}{4pt}
\begin{tabular}{@{}p{1.1cm}p{2.0cm}p{11.2cm}@{}}
\toprule
\textbf{Tahap} & \textbf{Sesi} & \textbf{Kegiatan} \\
\midrule
1 & 11 & Kickoff: pembagian 6 kelompok, pemilihan topik, pembagian peran \\
2 & 12 & Design awal (blok diagram), skematik \texttt{Fritzing}, 3D design
       enclosure, dan \textbf{Progress Report 1} --- \textbf{(Milestone,
       bobot 10\%)} \\
3 & 13 & Implementasi awal firmware + keamanan dasar (TLS, autentikasi);
       \textbf{Progress Report 2}: demo prototype awal \\
4 & 14 & Integrasi menyeluruh (sensor--ESP32--cloud--dashboard), debugging,
       dan uji coba \\
5 & 15 & Finalisasi laporan, rehearsal, dan \textbf{presentasi +
       demonstrasi} --- \textbf{(bobot 40\%)} \\
\bottomrule
\end{tabular}
\end{center}

Progress report disampaikan singkat setiap pertemuan sebagai bagian penilaian
proses (tercatat pada rubrik kerja tim).

\section{Keluaran (Deliverables)}

\begin{enumerate}
  \item \textbf{Prototype perangkat keras} yang berfungsi (bukan hanya
        simulasi Wokwi), dalam enclosure/rangkaian rapi.
  \item \textbf{Kode program} firmware dan sisi cloud/dashboard, diunggah ke
        repositori GitHub kelompok dengan README cara replikasi.
  \item \textbf{Desain teknis}: blok diagram, skematik Fritzing, file 3D
        enclosure (STL), dan daftar komponen beserta estimasi biaya.
  \item \textbf{Laporan} (5--10 halaman): latar belakang, desain sistem,
        implementasi, hasil pengujian (metrik), kendala, dan pembagian tugas
        anggota.
  \item \textbf{Video demo cadangan} (2--3 menit) untuk antisipasi kendala
        saat demo langsung.
  \item \textbf{Slide presentasi} untuk pertemuan 15.
\end{enumerate}

\section{Kriteria Penilaian Proyek}

\begin{center}
\begin{tabular}{>{\raggedright\arraybackslash}p{4.6cm}>{\raggedright\arraybackslash}p{7.9cm}p{1.8cm}}
\toprule
\textbf{Komponen} & \textbf{Deskripsi} & \textbf{Bobot} \\
\midrule
Fungsionalitas \& integrasi cloud & Alur end-to-end berfungsi: sensor
  $\rightarrow$ ESP32 $\rightarrow$ MQTT/REST $\rightarrow$ dashboard
  $\rightarrow$ kendali balik & 25\% \\
Desain \& implementasi hardware & Kualitas rangkaian/solder, kesesuaian
  skematik Fritzing, kerapian enclosure 3D & 15\% \\
Cakupan tingkat fitur & Tingkat 1 tuntas + minimal 2 pilihan Tingkat 2
  berfungsi; metrik keberhasilan tercapai & 15\% \\
Laporan \& dokumentasi & Kelengkapan laporan, README repositori, dokumentasi
  pengujian & 15\% \\
Presentasi \& demonstrasi & Kejelasan penyampaian, live demo lancar, kemampuan
  menjawab pertanyaan teknis & 15\% \\
Kerja tim \& progres & Kontribusi merata (progress report 1 \& 2, catatan
  bimbingan) & 10\% \\
Kreativitas & Fitur tambahan di luar spesifikasi yang relevan & 5\% \\
\midrule
\textbf{Total} & & \textbf{100\%} \\
\bottomrule
\end{tabular}
\end{center}

Nilai proyek akhir $=$ 10\% (milestone desain) $+$ 40\% (tabel di atas
$\times$ bobot proyek), sesuai struktur penilaian pada RPS.

\section{Platform dan Komponen}

\subsection*{Perangkat lunak (gratis/open source)}

\begin{itemize}
  \item \texttt{Wokwi} (\url{https://wokwi.com}) --- simulasi rangkaian
        \textbf{wajib sebelum perakitan}.
  \item \texttt{Fritzing} --- skematik dan breadboard view.
  \item \texttt{Tinkercad}/\texttt{Fusion 360} --- desain enclosure 3D.
  \item \texttt{Arduino IDE 2.0} --- firmware ESP32.
  \item Broker MQTT: \texttt{HiveMQ Cloud} (gratis) atau \texttt{Mosquitto};
        klien uji: \texttt{MQTT Explorer}.
  \item Dashboard: \texttt{Node-RED}, \texttt{Grafana}, atau halaman web
        mandiri; opsi \texttt{ThingSpeak}/\texttt{Blynk} free tier.
\end{itemize}

\subsection*{Estimasi biaya komponen per kelompok}

Komponen inti (ESP32, breadboard, kabel) dimanfaatkan dari praktikum
pertemuan 1--10. Komponen tambahan per topik diperkirakan Rp75.000--150.000
(sensor spesifik, relay, servo/pompa, OLED). Rencanakan belanja sesuai bill
of materials pada milestone desain.

\section{Ketentuan}

\begin{itemize}
  \item Setiap anggota wajib memahami seluruh bagian sistem, bukan hanya
        bagiannya; penilaian presentasi dapat menyasar siapa pun.
  \item \textbf{Aturan simulate-first}: rangkaian baru boleh dirakit pada
        hardware setelah lolos simulasi Wokwi dan disetujui dosen/asisten
        pada progress report 1.
  \item Layanan cloud digunakan pada \textbf{tier gratis}; dilarang memasukkan
        kartu kredit/membuat langganan berbayar atas nama institusi.
  \item Pembayaran (Topik 1) menggunakan \textbf{sandbox/mock}, bukan uang
        sungguhan.
  \item \textbf{Upaya dinilai}: tidak semua kelompok harus mencapai Tingkat
        3. Penilaian mengikuti struktur bertingkat --- Tingkat 1 yang berjalan
        sempurna, dokumentasi rapi, dan progress report jujur tentang apa yang
        berhasil/gagal tetap memperoleh nilai baik. Catat capaian tiap tingkat
        pada laporan.
  \item Penggunaan alat bantu AI generatif diperbolehkan untuk keperluan
        pembelajaran (menjelaskan konsep, debug), namun \textbf{wajib
        dideklarasikan} pada laporan; menyalin hasil mentah tanpa pemahaman
        akan terdeteksi pada sesi tanya jawab demo.
  \item Proyek yang terbukti menyalin karya kelompok lain dinilai 0 untuk
        seluruh anggota.
\end{itemize}

\end{document}
